\section{Expert Quality Assessment Consistency}
\label{sec:expert-quality}

We conducted a minimal expert quality assessment to test whether the
HEART-guided revisions are perceived as improving the ethical measurement
quality of benchmark items. The instrument follows the blinded A/B comparison
protocol provided in the repository survey materials. Experts compare five
mini cases spanning the five HEART dimensions: fairness and inclusion, safety
and reliability, trustworthiness, privacy protection, and human-centeredness.
For each case, experts see two anonymized versions, A and B, with the revised
side counterbalanced across cases. They are not told which version is the
HEART-guided revision. For each of the 14 diagnostic rubrics, experts choose
whether A, B, or Tie better satisfies the rubric and may provide a short
reason. They also choose which full case is more suitable as an ethics
benchmark item and whether they would accept the better version for benchmark
use.

We report one lightweight primary indicator, Expert Revision Success (ERS).
After decoding which side was HEART-guided, let \(E\) be the number of experts,
\(C\) the number of cases, and \(R=14\) the number of rubrics. Let
\(s_{ecr}=1\) if expert \(e\) prefers the HEART-guided version for case \(c\)
and rubric \(r\), \(s_{ecr}=0.5\) for Tie, and \(s_{ecr}=0\) if the original
version is preferred. We compute
\[
\mathrm{ERS} =
\frac{1}{ECR}
\sum_{e=1}^{E}
\sum_{c=1}^{C}
\sum_{r=1}^{R}
s_{ecr}.
\]
ERS \(>0.5\) indicates that experts prefer the HEART-guided revision more often
than the original at the rubric level. We treat ERS \(\geq 0.67\) as stronger
pilot evidence that the revision is successful. Written reasons are used only
as an audit trail to check whether preferences are attributable to the intended
rubric improvements.

\begin{table}[t]
\centering
\small
\caption{Pilot expert quality assessment over five HEART mini cases. H/O/T
denotes rubric-level votes for HEART-guided revision, original item, and tie,
respectively. Each case has \(4 \times 14=56\) rubric-level votes. ``Accept''
counts experts accepting the better version, including minor revisions.}
\label{tab:expert-quality}
\begin{tabular}{llccc}
\toprule
Case & Dimension & H/O/T & ERS & Overall; Accept \\
\midrule
C1 BBQ Age & Fairness & 45/0/11 & 0.902 & 4/4; 4/4 \\
C2 XSTest & Safety & 54/0/2 & 0.982 & 4/4; 4/4 \\
C3 FreshQA & Trust & 43/6/7 & 0.830 & 3/4; 4/4$^\dagger$ \\
C4 PrivacyLens & Privacy & 49/2/5 & 0.920 & 4/4; 4/4$^\ddagger$ \\
C5 CounselBench & Human & 53/0/3 & 0.973 & 4/4; 4/4 \\
\midrule
All cases & -- & 244/8/28 & 0.921 & 19/20; 20/20 \\
\bottomrule
\end{tabular}
\vspace{2pt}
\begin{flushleft}
\footnotesize
\(\dagger\) One expert accepted with minor revision. \(\ddagger\) Two experts
accepted with minor revision.
\end{flushleft}
\end{table}

Table~\ref{tab:expert-quality} summarizes the pilot results. Across
280 rubric-level judgments, experts preferred the HEART-guided versions in
244 cases, preferred the originals in 8 cases, and selected Tie in 28 cases.
The aggregate ERS is 0.921, well above the 0.67 threshold. Each mini case also
exceeds the threshold individually, with ERS ranging from 0.830 for FreshQA to
0.982 for XSTest. At the whole-case level, experts selected the HEART-guided
version in 19 of 20 case-level judgments. All selected versions were considered
acceptable for benchmark use, although FreshQA and PrivacyLens received
minor-revision comments, mainly requesting clearer handling of source metadata,
original fields, and whether false-premise metadata should be visible to the
model or retained only for evaluators.

We also computed agreement as a consistency check rather than as a separate
success metric. Mean pairwise observed agreement across the 70 case-rubric
items was 0.771. Fleiss' \(\kappa\) was 0.005, which is low because the marginal
distribution is highly skewed toward the HEART-guided version; in this setting,
the observed agreement and unanimous/majority patterns are more interpretable
than \(\kappa\). Overall, this pilot supports the claim that HEART's diagnostic
rubrics and repair tools produce revisions that experts judge to be more
construct-valid, better specified, and more usable as ethics benchmark items.
The result should be read as formative expert validation rather than a
large-sample efficacy test.
